\documentclass{article} % For LaTeX2e
\usepackage{iclr2025_conference,times}

\usepackage{hyperref}
\usepackage{url}
\usepackage{booktabs}

\title{Live Tennis Win Probability from Score-State Recursions and In-Match Serve Updating}

% Authors must not appear in the submitted version. They should be hidden
% as long as the \iclrfinalcopy macro remains commented out below.
\author{Anonymous Author(s)}

\newcommand{\fix}{\marginpar{FIX}}
\newcommand{\new}{\marginpar{NEW}}
\newcommand{\pserve}{p_{\mathrm{serve}}}

%\iclrfinalcopy % Uncomment for camera-ready version, but NOT for submission.
\begin{document}

\maketitle

\begin{abstract}
Live win probability in tennis is a useful test case for combining exact game
structure with learned player information. The score evolves through nested
point, game, set, and match states, but the players' pre-match strengths and
same-day serving form are only partially observed. I build a live tennis win
probability model on Grand Slam point-by-point data from 2011--2024. A
hierarchical Markov recursion converts two player-specific serve-point win
probabilities and the current score into an exact match win probability.
Pre-match Elo ratings set the prior serve advantage, and in-match serve results
update that prior through a simple shrinkage estimator. On held-out 2023--2024
Grand Slam matches, the symmetric score-only recursion reaches 0.5363 log loss
at the midpoint of matches, the Elo-asymmetric recursion improves to 0.4549, and
the serve-shrink model evaluated at 70\% match progress reaches 0.3000 log loss
before a light logistic calibration layer improves it to 0.2940. The result is a
small, inspectable model whose largest gains come from giving the structural
tennis recursion the right player-specific inputs.
\end{abstract}

\section{Introduction}

Tennis has an unusually clean live prediction structure. At any instant, the
score determines how many points, games, sets, and matches remain; the rules
then determine exactly how point-level advantage propagates upward. This makes a
pure machine-learning framing somewhat wasteful. A flexible model can learn that
a player two sets up is likely to win, but the sport already tells us how
valuable that lead is under any assumed point-winning probabilities.

The harder part is choosing those point-winning probabilities. A score-only
model treats the two players as exchangeable, so every match begins at 50--50.
That is a poor description of a first-round match between a top seed and a
qualifier. Conversely, a pre-match model knows the players' long-run strength
but not whether one of them is serving especially well today. The model in this
paper keeps the rule-based recursion and focuses the learning problem on the
quantities the recursion actually needs: the chance that each player wins a
point on serve.

The contribution is deliberately modest and interpretable. I compare three
families of live predictors: a symmetric Markov anchor using the score alone, an
asymmetric Markov model whose serve probabilities are shifted by pre-match Elo
difference, and a serve-shrink model that blends the Elo prior with serve points
observed earlier in the same match. The approach uses only information available
live, and each fitted parameter has a direct tennis interpretation.

\section{Data}

The point-level data comes from Jeff Sackmann's public Grand Slam
point-by-point files, covering the Australian Open, Roland Garros, Wimbledon,
and the US Open from 2011 through 2024. After filtering incomplete or unscored
matches, the prepared table contains 1,505,355 points from 8,222 matches. Each
row is oriented to the file's player 1 and is labeled by whether player 1
eventually won the match.

For each point I store the score after the point: sets won, games won in the
current set, point score inside the current game or tiebreak, server, best-of
length, and whether the set is in a tiebreak. I also compute live context
features using only earlier points in the match: each player's serve points
played, serve points won, aces, ace rates, and rally-length summaries. The
running serve rates default to 0.5 before a player has served a point, avoiding
any lookahead from the current point.

Pre-match player strength is supplied by an Elo table built from ATP and WTA
tour-level match results from 2005 onward. Ratings are updated chronologically
before each match, with a decaying step size as a player accumulates matches.
The resulting pre-match ratings are joined to Grand Slam point-by-point matches
by year, event, round ordering, and normalized player names. The Elo artifact
matches 7,896 of the 8,222 prepared Grand Slam matches.

All tuning uses a time split. Matches through 2021 form the training set, 2022
is used for validation and parameter selection, and 2023--2024 are held out for
testing. The held-out period contains 179,832 points from 975 matches in the
raw point table. Because the current scripts evaluate one live state per match,
fixed-fraction experiments contain 969 held-out matches after the Elo join.

\section{Method}

\subsection{Markov score recursion}

Let $p_a$ be the probability that player 1 wins a point on player 1's serve,
and let $p_b$ be the probability that player 2 wins a point on player 2's
serve. Given $(p_a,p_b)$, the model computes player 1's match win probability
from the current score by a nested recursion.

At the game level, $G(s,i,j)$ is the probability that the server wins the game
from point score $i$--$j$ when the server wins each point with probability $s$.
Absorbing states occur once either player has at least four points and leads by
two. At deuce and advantage states, the implementation uses the closed-form
probability of eventually winning the deuce cycle. Set probabilities are then
computed from game probabilities, alternating server after each game. At
6--6, a tiebreak recursion handles the one-point, then two-point serve
alternation. Finally, a match recursion combines the probability of winning the
current set with the probability of winning future fresh sets until one player
reaches the required number of sets.

The recursion is memoized over rounded serve probabilities and score states, so
the model can score a large point table while solving each distinct state only
once. This structure enforces tennis rules exactly while leaving player quality
to the choice of $p_a$ and $p_b$.

\subsection{Symmetric anchor}

The first baseline assumes equal players:
\[
    p_a = p_b = p.
\]
The only free parameter is the global serve-point win probability $p$, selected
on 2022 validation states from the grid $\{0.60,0.61,\ldots,0.65\}$. The
selected value is $p=0.64$. This model is a pure score-state anchor: two
matches with the same score receive the same probability, regardless of player
identity or match history.

\subsection{Elo-asymmetric serve probabilities}

The asymmetric model keeps the same recursion but assigns different serve
probabilities using the pre-match Elo gap:
\[
    e = \mathrm{clip}(\alpha(\mathrm{Elo}_1-\mathrm{Elo}_2), -0.15, 0.15),
    \quad
    p_a = \mathrm{clip}(\beta + e, 0.45, 0.88),
    \quad
    p_b = \mathrm{clip}(\beta - e, 0.45, 0.88).
\]
The validation grid selects $\beta=0.61$ and $\alpha=1.3\times 10^{-4}$. This
single change lets the pre-match prediction start away from 50--50 while
preserving the exact score dynamics.

\subsection{Serve-shrink updating}

The full structural model updates each player's serve probability with
same-match evidence. Let $r_1,n_1$ be player 1's serve-point win rate and serve
points observed before the current row, and define $r_2,n_2$ analogously for
player 2. With Elo-implied priors $\pi_1=\beta+e$ and $\pi_2=\beta-e$, the
serve probabilities are
\[
    p_a = \frac{n_1 r_1 + \kappa \pi_1}{n_1+\kappa},
    \qquad
    p_b = \frac{n_2 r_2 + \kappa \pi_2}{n_2+\kappa}.
\]
The pseudo-count $\kappa$ controls how quickly the model trusts in-match serve
performance. Validation over $\{40,80,160,320,640\}$ selects $\kappa=40$ in the
current run. This is still a conservative update: after 20 service points, the
prior receives twice as much weight as the observed serve rate.

I also fit an optional logistic calibration layer on top of the structural
prediction. Its input is the logit of the serve-shrink Markov probability plus
the live context features. A median imputer, standard scaler, and
$L_2$-regularized logistic regression are fit on the training and validation
seasons after selecting the regularization strength on 2022. This layer is not
needed to define the structural model, but it provides a small empirical
calibration improvement and a way to test whether additional live features add
signal beyond the Markov score.

\section{Experiments}

The current model scripts evaluate a single live state per match at a fixed
fraction of match progress. The symmetric and Elo-asymmetric Markov models are
evaluated at the midpoint of each match; the serve-shrink model is evaluated at
70\% of match progress, where in-match serve evidence is more informative. The
gradient-boosted score baseline is implemented with XGBoost over raw score
and live context features.

\begin{table}[t]
\caption{Held-out results on 2023--2024 Grand Slam matches. Each row reports
the corresponding script's current fixed-fraction evaluation.}
\label{tab:results}
\begin{center}
\begin{tabular}{lcccc}
\toprule
Model & Match fraction & Log loss & Brier & Accuracy \\
\midrule
Symmetric Markov & 0.50 & 0.5363 & 0.1745 & 0.7703 \\
Elo-asymmetric Markov & 0.50 & 0.4549 & 0.1459 & 0.8050 \\
Serve-shrink Markov & 0.70 & 0.3000 & 0.0972 & 0.8607 \\
Serve-shrink + logistic calibration & 0.70 & 0.2940 & 0.0959 & 0.8576 \\
\bottomrule
\end{tabular}
\end{center}
\end{table}

Table~\ref{tab:results} shows the main pattern. The score-only recursion is
already strong because tennis scoring is highly informative: at the midpoint of
a match, the current set and game state explain much of the outcome. Adding the
Elo prior produces the largest fair head-to-head improvement, reducing midpoint
log loss from 0.5363 to 0.4549 and raising accuracy from 0.7703 to 0.8050. This
gain comes from fixing the otherwise unrealistic assumption that the two players
are equally strong before any points are played.

The serve-shrink model is evaluated later in matches, so its numbers should not
be read as a direct apples-to-apples comparison with the midpoint models.
Instead, they show the intended live behavior: once enough match-specific serve
data has accumulated, updating the point-level inputs to the recursion yields a
substantially sharper forecast. The calibration layer lowers log loss further
from 0.3000 to 0.2940, although its 0.5-threshold accuracy is slightly lower,
which is consistent with a probability calibration improvement rather than a
hard-classification improvement.

\subsection{Feature importance}

Permutation importance for the calibrated serve-shrink model confirms that the
model is driven by tennis score structure first and live serving form second.
On the 70\% held-out states, shuffling player 1 sets increases log loss by
0.4835, shuffling player 2 sets by 0.4679, and shuffling best-of length by
0.3072. The largest non-score features are the observed serve rates:
player 2's serve rate increases log loss by 0.1494 when permuted, and player
1's by 0.1473. Elo difference remains important, increasing log loss by 0.0474,
while rally-length features have essentially zero marginal importance in the
current linear calibration layer.

\section{Discussion}

The results support a simple modeling lesson: in a sport with known recursive
rules, the most useful learning often happens at the interface between
observations and the structural model. The Markov recursion does not need to
learn that winning a tiebreak changes a set, or that a two-set lead in a
best-of-five match is valuable. Those facts are exact. What it needs are
reasonable player-specific point probabilities.

The Elo-asymmetric model gives the recursion a sensible pre-match prior and
accounts for most of the improvement over the symmetric anchor. The serve-shrink
model then handles the live part of the problem: today's serving performance can
depart from long-run strength, but it should be shrunk heavily early in the
match. The selected pseudo-count of 40 serve points is interpretable as the
amount of prior evidence attached to the Elo-implied serve probability.

There are also clear limitations. The current evaluation scripts use fixed
match fractions rather than reporting every point on one common test table, so
the headline rows answer slightly different live questions. The point
probability model is also intentionally narrow: return strength, surface,
fatigue, injury, pressure, and tactical changes all enter only indirectly
through Elo or observed serve outcomes. Finally, the Grand Slam join to Elo is
name-based and misses a small share of matches, which may introduce selection
effects if the unmatched matches are not random.

\section{Conclusion}

This paper drafts a compact live tennis win-probability system around an exact
score recursion. A symmetric Markov model provides a transparent score-only
anchor, Elo-derived asymmetric serve probabilities supply the major pre-match
strength correction, and Bayesian-style serve shrinkage updates the model with
same-day evidence. The resulting model is small enough to inspect by hand, fast
enough to run point by point, and accurate enough to serve as a credible
baseline for richer live tennis models.

\subsubsection*{Acknowledgments}
The data preparation relies on the public tennis datasets maintained by Jeff
Sackmann. Any remaining errors in data joining, modeling, or interpretation are
my own.

\begin{thebibliography}{9}

\bibitem[Sackmann(2024a)]{sackmann-slam}
Jeff Sackmann.
\newblock Grand Slam point-by-point tennis data.
\newblock \url{https://github.com/JeffSackmann/tennis_slam_pointbypoint}.

\bibitem[Sackmann(2024b)]{sackmann-atp}
Jeff Sackmann.
\newblock ATP tennis rankings, results, and stats.
\newblock \url{https://github.com/JeffSackmann/tennis_atp}.

\bibitem[Sackmann(2024c)]{sackmann-wta}
Jeff Sackmann.
\newblock WTA tennis rankings, results, and stats.
\newblock \url{https://github.com/JeffSackmann/tennis_wta}.

\end{thebibliography}

\appendix
\section{Implementation Details}

The implementation is organized as a small pipeline. \texttt{prepare\_data.py}
builds the point-state table and live running features.
\texttt{elo.py} builds pre-match ratings from tour-level match history and
joins them to the point table. \texttt{markov.py} contains the memoized
point-game-set-match recursion used by all structural models.
\texttt{baseline\_markov.py}, \texttt{asymmetric\_markov.py}, and
\texttt{serve\_shrink\_model.py} run the three main experiments. The calibrated
serve-shrink model and permutation importances are saved in the
\texttt{artifacts} directory.

\end{document}
